% AUTO-GENERATED by scripts/build_topics.py — do not edit.
% Edit topics/bodies/topic_34.tex or topics/preamble.tex instead.
%% Placeholder. Replace with the written chapter.
\textit{Източници:} PDF \texttt{33.pdf}, снимки \texttt{34та тема}

\subsection*{Анотация от конспекта}
Итерационни методи за решаване на нелинейни уравнения. Да се дефинира понятието неподвижна точка на изображението  и да се докаже, че ако  е непрекъснато изображение на интервала [a, b] в себе си, то  има поне една неподвижна точка в [a, b]. Да се покаже, че решаването на уравнението f(x) = 0 може да се сведе към намиране на неподвижна точка. Да се дефинира понятието свиващо изображение и да се докаже, че ако  е непрекъснато изображение на интервала [a, b] в себе си и е свиващо с константа на Липшиц q< 1, то: а) уравнението x= (x) има единствен корен  в [a, b]; б) редицата \{xn\} от последователни приближения (при произволно x0[a, b] и xn+1 = (xn), n = 0, 1, 2, ..., клони към  при n→, като |xn – |  (b – a)qn, за всяко n. Да се получи като следствие, че ако  е корен на уравнението x= (x) и  има непрекъсната производна в околност U на , за която |′(ξ)| < 1, то при достатъчно добро начално приближение x0 итерационният процес, породен от , е сходящ със скоростта на геометрична прогресия. Да се дефинира понятието ред на сходимост. Да се дадат геометрична илюстрация, формула за последователните приближения и ред на сходимост при: метод на хордите, метод на секущите и метод на Нютон. Да се докаже, че при метода на хордите сходимостта е със скоростта на геометричната прогресия (при условие, че коренът е отделен в достатъчно малък интервал). Литература: [1], [2], [16].
